\documentclass[10pt,a4paper]{amsart}
\usepackage[margin=19mm]{geometry}
\usepackage{amsmath,amssymb,mathtools}
\usepackage[scr=rsfs,cal=euler]{mathalfa}
\usepackage{xfrac}
\usepackage[unicode,hidelinks,bookmarks=false]{hyperref}
\newcommand{\ie}{i.e.\ }
\newcommand{\cf}{cf.\ }
\newcommand{\resp}{respectively\ }
\newcommand{\Subsection}{\subsection}
\providecommand{\nobreakdash}{\nobreak\hbox{-}}
\setlength{\emergencystretch}{2em}
\multlinegap0pt
\usepackage{amssymb}
\usepackage{algorithm}\usepackage{algpseudocode}
\usepackage{float}
\newtheorem{theorem}{Theorem}
\title{Vehicle Routing Problem with Resource-Constrained Pickup and Delivery}
\author{Manbir Sodhi, Romesh Prasad, Harishjitu Saseendran}
\date{Source: arXiv:2602.23685v2 (CC BY 4.0)}
\begin{document}
\maketitle

\noindent\emph{Source and licence.} Vehicle Routing Problem with Resource-Constrained Pickup and Delivery. Version arXiv:2602.23685v2 (CC BY 4.0), distributed by arXiv under the Creative Commons Attribution 4.0 International licence, http://creativecommons.org/licenses/by/4.0/, as recorded in the arXiv licence metadata for this exact version and shown on the abstract page https://arxiv.org/abs/2602.23685v2. Full licence notice: \texttt{licenses/arxiv-2602.23685-CC-BY-4.0.txt}.\par\smallskip

\noindent\textit{Editorial scope.} Counted unit: the article's theorem theorem (source0.tex lines 288--290). The article's complete original proof of it is delivered with it (source0.tex lines 292--303, one \texttt{proof} environment of 12 lines). No mathematics is written, repaired or re-derived: the statement and the proof are the article's own lines, in the article's own notation, and no retained line is substituted or re-flowed.  No further source-local context is needed: every cross-reference of every retained line resolves inside the two delivered spans.  Nothing was omitted from any retained span, so the package carries no omission record.  The retained lines cite 1 works, whose entries are transcribed verbatim from the article's own bibliography source in the e-print and delivered with short editorial labels recorded in the item's reference mapping.  No retained line contains a figure environment, an image-inclusion command or drawing code, so no image is delivered and the package declares no asset dependency.  Editorial formatting only, in addition: an AMS \texttt{amsart} preamble that restates the article's own declarations and macros used by the retained lines, namely the environments \texttt{theorem} on separate counters, together with the standard packages those lines need.  The retained environments print in this document as Theorem 1 in order of appearance, whereas the article prints the same items with the numbering of its own full document.  The complete native source of the article as distributed in the arXiv e-print for this exact version is bundled unchanged as \texttt{sources/195416/source0.tex}.
\begin{theorem}
VRP-RPD is NP-hard.
\end{theorem}

\begin{proof}
By reduction from the min-max Vehicle Routing Problem (min-max VRP). 
Consider a VRP-RPD instance where processing times $p_c = 0$ for all 
customers $c \in C$, and vehicle capacity $k \geq n$. Since processing 
time is zero, pickup can occur immediately after dropoff at each customer. 
Since capacity is non-binding, each customer requires exactly one visit 
where both operations are performed instantaneously by the same vehicle. 
The problem reduces to partitioning customers among $m$ vehicles and 
sequencing visits to minimize the maximum return time, precisely the 
min-max VRP. Since min-max VRP is NP-hard~\cite{lenstra1981complexity}, 
VRP-RPD is NP-hard.
\end{proof}

\begin{thebibliography}{99}
\bibitem[lenstr]{lenstra1981complexity} J.K. Lenstra and A.H.G. Rinnooy~Kan. Complexity of vehicle routing and scheduling problems. {\em Networks}, 11(2):221--227, 1981.
\end{thebibliography}
\end{document}
